%==============================================================================
%  Preámbulo del informe: paquetes y configuración.
%  Lo comparten Plantilla_Informe_TFG.tex y Apendice1_Implicaciones.tex, para que
%  la copia aparte del Apéndice 1 tenga exactamente el mismo formato.
%==============================================================================

%------------------------------------------------------------------ Idioma ---
\usepackage[T1]{fontenc}
\usepackage[utf8]{inputenc}
\usepackage{lmodern}
% es-tabla: fuerza «Tabla» en lugar de «Cuadro» (predeterminado de babel-spanish).
% es-noquoting: evita que babel se apropie de los signos " y < >.
\usepackage[spanish,es-tabla,es-noquoting]{babel}

%----------------------------------------------------- Página y espaciado ----
\usepackage[letterpaper,margin=1in]{geometry}
\usepackage{setspace}
\onehalfspacing                 % interlineado 1,5
\usepackage{parskip}            % separación entre párrafos, sin sangría

%------------------------------------------------------- Pie de página -------
\usepackage{fancyhdr}
\pagestyle{fancy}
\fancyhf{}
\renewcommand{\headrulewidth}{0pt}
\fancyfoot[R]{\thepage}         % número de página a la derecha (como el .docx)
\fancypagestyle{plain}{\fancyhf{}\renewcommand{\headrulewidth}{0pt}\fancyfoot[R]{\thepage}}

%------------------------------------------------- Tablas, figuras, código ---
\usepackage{graphicx}
\graphicspath{{media/}{figuras/}}
\usepackage{booktabs}
\usepackage{tabularx}
\usepackage{longtable}
\usepackage{array}
\usepackage{float}
\usepackage{enumitem}
\setlist{nosep,topsep=4pt}

\usepackage{caption}
\captionsetup{font=small,labelfont=bf,labelsep=colon,skip=6pt}
\captionsetup[table]{position=above}
\captionsetup[figure]{position=below}

% Columnas de ancho elástico para tabularx:
%   C = centrada,  L = alineada a la izquierda
% Anteponga \ancho{f} para darle a una columna f veces el ancho base
% (la suma de los factores debe ser igual al número de columnas X).
\newcolumntype{C}{>{\centering\arraybackslash}X}
\newcolumntype{L}{>{\raggedright\arraybackslash}X}
\newcommand{\ancho}[1]{\hsize=#1\hsize\linewidth=\hsize}

% Lista para usar DENTRO de una celda de tabla: sin espacio superior, para que
% la primera viñeta quede alineada con el texto de las demás columnas.
\newlist{celdaitem}{itemize}{1}
\setlist[celdaitem]{label=\textbullet,leftmargin=1.2em,nosep,topsep=0pt,
                    partopsep=0pt,parsep=0pt,itemsep=2pt}

% Un poco más de aire entre las filas de las tablas
\renewcommand{\arraystretch}{1.15}

%------------------------------------------------------------- Títulos -------
\usepackage{titlesec}
\titleformat{\section}{\normalfont\Large\bfseries}{\thesection.}{0.6em}{}
\titleformat{\subsection}{\normalfont\large\bfseries}{\thesubsection.}{0.6em}{}
\titleformat{\subsubsection}{\normalfont\normalsize\bfseries}{\thesubsubsection.}{0.6em}{}
\titlespacing*{\section}{0pt}{18pt}{10pt}
\titlespacing*{\subsection}{0pt}{14pt}{8pt}
\titlespacing*{\subsubsection}{0pt}{12pt}{6pt}
\setcounter{secnumdepth}{3}     % numeración hasta 1.1.1
\setcounter{tocdepth}{3}        % índice general hasta 1.1.1

%--------------------------------------------------------- Hiperenlaces ------
\usepackage[hidelinks,unicode]{hyperref}
\hypersetup{
  pdftitle   = {Generalización de framework de selección automática de pipelines de
                preprocesamiento a múltiples tareas de aprendizaje automático mediante
                meta-aprendizaje},
  pdfauthor  = {Ricardo Antonio Borbón Mena},
  pdfsubject = {Informe de Trabajo Final de Graduación},
  pdflang    = {es-CR},
  bookmarksnumbered = true
}
\usepackage{bookmark}

%-------------------------------------------- Comando para el texto guía -----
% Todo el texto de ayuda de la plantilla va dentro de \guia{...}. Aparece en
% gris y cursiva para distinguirlo. BÓRRELO conforme escriba su informe; al
% terminar no debe quedar ningún \guia en el documento.
\usepackage{xcolor}
\newcommand{\guia}[1]{{\itshape\color{gray!75!black}#1}}

% Atajo para escribir nombres de comandos dentro del texto guía
\newcommand{\cmd}[1]{\texttt{\textbackslash #1}}

%------------------------------------------------------- Bibliografía --------
% ESTÁNDAR DE CITACIÓN: IEEE  (decidido para este TFG; los lineamientos exigen
% usar uno solo — APA 7 o IEEE — de forma rigurosa en TODO el documento).
%
% Las referencias se administran en el archivo referencias.bib.
% En el texto se cita con \cite{clave}, que produce marcas numéricas [1], [2]…
% Requiere compilar con biber (latexmk lo hace automáticamente).
\usepackage{csquotes}   % requerido por biblatex cuando se usa babel
\usepackage[backend=biber,style=ieee,sorting=none]{biblatex}
\addbibresource{referencias.bib}

% sorting=none: las referencias se numeran en el orden en que aparecen citadas
% en el texto, como manda IEEE.

